\documentclass[11pt,a4paper]{article}
\usepackage[utf8]{inputenc}
\usepackage[T1]{fontenc}
\usepackage{amsmath,amssymb,amsthm,geometry,booktabs,array,hyperref,mathtools}
\usepackage{xcolor,colortbl,setspace,enumitem,fancyhdr,longtable,tikz}

\geometry{margin=2.6cm, top=3.2cm, bottom=3.2cm}
\onehalfspacing

\definecolor{deepblue}{RGB}{10,60,120}
\definecolor{lightgray}{gray}{0.94}
\definecolor{levelA}{RGB}{220,235,255}
\definecolor{levelB}{RGB}{240,250,235}

\hypersetup{colorlinks=true,linkcolor=deepblue,citecolor=deepblue,urlcolor=deepblue,
  pdftitle={QGD Chapter 5: Hamiltonian Formalism},pdfauthor={Romeo Matshaba}}

\pagestyle{fancy}
\fancyhf{}
\fancyhead[L]{\small\textit{QGD Chapter 5}}
\fancyhead[R]{\small\textit{Hamiltonian Formalism}}
\fancyfoot[C]{\thepage}
\renewcommand{\headrulewidth}{0.4pt}

\theoremstyle{plain}
\newtheorem{theorem}{Theorem}[section]
\newtheorem{lemma}[theorem]{Lemma}
\newtheorem{corollary}[theorem]{Corollary}
\newtheorem{proposition}[theorem]{Proposition}
\theoremstyle{definition}
\newtheorem{definition}[theorem]{Definition}
\theoremstyle{remark}
\newtheorem{remark}[theorem]{Remark}

\newcommand{\lQ}{\ell_Q}
\newcommand{\mQ}{m_Q}
\newcommand{\boxg}{\square_g}
\newcommand{\half}{\tfrac{1}{2}}
\newcommand{\Hperp}{H_\perp}
\newcommand{\Hi}{H_i}
\newcommand{\HADM}{H_{\mathrm{ADM}}}
\newcommand{\HEOB}{H_{\mathrm{EOB}}}
\newcommand{\HPU}{H_{\mathrm{PU}}}
\newcommand{\sigt}{\sigma_t}

\title{%
  \vspace{-1.2cm}
  {\large\scshape Quantum Gravitational Dynamics}\\[6pt]
  {\LARGE\bfseries Chapter 5: Hamiltonian Formalism}\\[6pt]
  {\large From the ADM Decomposition to the QGD EOB Hamiltonian,\\
  Pais--Uhlenbeck Structure, and Connection to IMRI Waveforms}
}
\author{Romeo Matshaba\\[3pt]
  \small Department of Physics, UNISA\quad
  DOI: \href{https://doi.org/10.5281/zenodo.18827993}{10.5281/zenodo.18827993}}
\date{March 2026}

\begin{document}
\maketitle\thispagestyle{empty}

\begin{abstract}
\noindent
We develop the complete Hamiltonian formalism for Quantum Gravitational Dynamics,
proceeding in six levels from the covariant master action to observable gravitational
wave frequencies.  Starting from the 3+1 ADM decomposition of the QGD $\sigma$-field
action, we derive the QGD Hamiltonian constraint (which differs from the GR constraint
by a $\sigma$-field wave term), the canonical Poisson structure, and the
Pais--Uhlenbeck factorisation into massless and Planck-mass sectors.  The reduced
two-body Hamiltonian is shown to reproduce Newton's law from the $M$-matrix cross-term,
and the effective one-body (EOB) Hamiltonian is identified with the $A(u;\nu)$ potential
of Chapter~4.  We then establish the connection between the QGD dual-ladder structure
and the hybrid PN+BHP approach of Laxman, Fujita, and Mishra~\cite{LFM2026} for
eccentric IMRIs: Ladder~A corresponds to the PN near-zone sector, Ladder~B to the
BHP/tail sector.  QGD provides a \emph{unified, parameter-free} formulation that
subsumes both sectors from a single Lagrangian.
\end{abstract}

\tableofcontents\newpage

%=============================================================
\section{The GR Hamiltonian Landscape}
\label{sec:GR_Ham}
%=============================================================

Before deriving the QGD Hamiltonian, we survey the three GR Hamiltonian
frameworks that QGD generalises.

\subsection{ADM Hamiltonian}

The ADM formalism (Arnowitt--Deser--Misner, 1959--1962) provides a Hamiltonian
formulation of GR by decomposing four-dimensional spacetime into a foliation
of three-dimensional spatial hypersurfaces evolving in time.
The canonical variables are the 3-metric $\gamma_{ij}$ on each hypersurface $\Sigma_t$
and its conjugate momentum $\pi^{ij}$.

The 3+1 decomposition gives:
\begin{equation}
  ds^2 = -(N^2 - N_iN^i)\,dt^2 + 2N_i\,dx^i\,dt + \gamma_{ij}\,dx^i\,dx^j,
  \label{eq:ADM_metric}
\end{equation}
where $N$ is the lapse and $N^i$ is the shift.  The ADM action is
\begin{equation}
  S_{\mathrm{ADM}} = \int dt\,d^3x\,\Bigl[\pi^{ij}\dot\gamma_{ij}
  - N\Hperp^{\mathrm{GR}} - N^i\Hi^{\mathrm{GR}}\Bigr],
  \label{eq:ADM_action}
\end{equation}
with two sets of primary constraints: the scalar (Hamiltonian) constraint
$\Hperp^{\mathrm{GR}} = 0$ and the vector (momentum) constraint $\Hi^{\mathrm{GR}} = 0$.
The total Hamiltonian $\HADM = \int(N\Hperp + N^i\Hi)\,d^3x$ is identically zero
on solutions --- a problem that plagued GR for many years before the ADM analysis
revealed the 2 physical degrees of freedom per spatial point.

Hamiltonian formalisms provide powerful tools for the computation of approximate
analytic solutions of the Einstein field equations, particularly for the post-Newtonian
dynamics of compact binaries.

\subsection{EOB Hamiltonian}

The EOB Hamiltonian (Buonanno--Damour 1999) reduces the two-body problem to
one body in an effective geometry:
\begin{equation}
  H_{\mathrm{real}} = Mc^2\sqrt{1 + 2\nu\!\left(\frac{H_{\mathrm{eff}}}{\mu c^2}-1\right)},
  \label{eq:H_real}
\end{equation}
where
\begin{equation}
  H_{\mathrm{eff}}^2 = A(r;\nu)\!\left[\mu^2c^4
    + \frac{p_\phi^2 c^2}{r^2}
    + \frac{A(r;\nu)}{D(r;\nu)}p_r^2 c^2
    + Q(r;\nu,p_r)\right].
  \label{eq:H_eff}
\end{equation}
The potential $A(u;\nu)$ is exactly what we derived in Chapter~4.

\subsection{Pais--Uhlenbeck oscillator}

For a Lagrangian $L \propto q(\partial_t^2 + \omega_1^2)(\partial_t^2 + \omega_2^2)q$,
Ostrogradski's theorem gives a Hamiltonian with ghost modes.  The QGD master equation
has exactly this structure at the field-theory level; Section~\ref{sec:PU} shows
how ghost-freedom is maintained.

%=============================================================
\section{Canonical Variables for the QGD $\sigma$-Field}
\label{sec:canonical}
%=============================================================

\begin{definition}[QGD canonical variables]\label{def:canonical}
The QGD fundamental field $\sigma_\mu$ and its canonical momenta are:
\begin{align}
  \text{Coordinate: }&\sigma_\mu(x), \quad \mu = t,r,\theta,\phi,
  \label{eq:coord}\\
  \text{Momentum: }&\Pi^\mu(x) = \frac{\partial\mathcal{L}}{\partial(\partial_t\sigma_\mu)}.
  \label{eq:mom}
\end{align}
\end{definition}

For the kinetic Lagrangian $\mathcal{L}_{\mathrm{kin}} = (\hbar^2/2M)\sqrt{-g}\,g^{\mu\nu}\nabla_\mu\sigma^\alpha\nabla_\nu\sigma_\alpha$, the canonical momentum is
\begin{equation}
  \Pi^t = \frac{\hbar^2}{M}\sqrt{-g}\,g^{tt}\,\dot\sigma_t
        = -\frac{\hbar^2}{M}\cdot\frac{r^2}{\sqrt{f}}\,\dot\sigma_t,
  \label{eq:Pi_t}
\end{equation}
where $f = 1-\sigma_t^2$.  In isotropic gauge ($\sigma_r = 0$):
$\Pi^r = \Pi^\theta = \Pi^\phi = 0$ (constraint).

\begin{theorem}[Fundamental Poisson bracket]
\begin{equation}
  \{\sigma_\mu(t,\mathbf{x}),\,\Pi^\nu(t,\mathbf{x}')\}
  = \delta^\nu_\mu\,\delta^3(\mathbf{x}-\mathbf{x}').
  \label{eq:PB}
\end{equation}
\end{theorem}

%=============================================================
\section{3+1 Decomposition of the QGD Action}
\label{sec:3plus1}
%=============================================================

Applying the ADM foliation $\mathcal{M} = \mathbb{R}\times\Sigma$ to the QGD action:

\begin{theorem}[QGD in 3+1 form]\label{thm:3plus1}
The QGD master action decomposes as
\begin{equation}
  S_{\mathrm{QGD}} = \int\!dt\!\int_\Sigma\!d^3x
    \Bigl[\Pi^\mu\dot\sigma_\mu - N\,\Hperp^{\mathrm{QGD}} - N^i\Hi^{\mathrm{QGD}}\Bigr].
  \label{eq:S_3plus1}
\end{equation}
\end{theorem}

\subsection{Identification of the lapse and shift}

From the QGD metric with $\sigma_\mu = (\sigt,0,0,0)$:
\begin{equation}
  N = \sqrt{f} = \sqrt{1-\sigt^2}, \qquad N^i = 0, \qquad
  \gamma_{ij} = \mathrm{diag}(1,\,r^2,\,r^2\sin^2\theta).
  \label{eq:lapse_shift}
\end{equation}
The lapse $N = \sqrt{f}$ vanishes at the QGD horizon $\sigt = 1$: time ``stops''
at the horizon, exactly as in GR.

\subsection{The extrinsic curvature}

\begin{equation}
  K_{ij} = -\frac{1}{2N}\!\left(\dot\gamma_{ij} - D_iN_j - D_jN_i\right) = 0
  \label{eq:K}
\end{equation}
for the static Schwarzschild background ($\dot\gamma_{ij} = 0$, $N^i = 0$).

%=============================================================
\section{The QGD Hamiltonian Constraint}
\label{sec:constraint}
%=============================================================

\begin{theorem}[QGD Hamiltonian constraint]\label{thm:constraint}
The QGD scalar constraint takes the form
\begin{equation}
  \boxed{\Hperp^{\mathrm{QGD}} = \Hperp^{\mathrm{GR}}[\gamma,\pi]
    + H_\perp^{\mathrm{kin}}[\sigma,\Pi]
    + H_\perp^{\mathrm{grad}}[\sigma]
    + H_\perp^{\mathrm{self}}[\sigma]
    + H_\perp^{\mathrm{stiff}}[\sigma] = 0}
  \label{eq:constraint}
\end{equation}
with the five contributions:
\begin{align}
  \Hperp^{\mathrm{GR}}
  &= \frac{1}{\sqrt\gamma}\!\left[-\frac{1}{16\pi G}
      \left(\pi^{ij}\pi_{ij} - \tfrac{1}{2}\pi^2\right)
      + \frac{\sqrt\gamma\,{}^3\!R}{16\pi G}\right],
  \label{eq:H_GR}\\[4pt]
  H_\perp^{\mathrm{kin}}
  &= \frac{M}{2\hbar^2}\frac{(\Pi^t)^2}{\sqrt\gamma},
  \label{eq:H_kin}\\[4pt]
  H_\perp^{\mathrm{grad}}
  &= \frac{\hbar^2}{2M}\sqrt\gamma\,\gamma^{ij}D_i\sigt D_j\sigt,
  \label{eq:H_grad}\\[4pt]
  H_\perp^{\mathrm{self}}
  &= \frac{\hbar^2}{2M}\sqrt\gamma\cdot
    \frac{\sigt^2}{f}(\nabla\sigt)^2,
  \label{eq:H_self}\\[4pt]
  H_\perp^{\mathrm{stiff}}
  &= \frac{\kappa\lQ^2\hbar^2}{2M}\sqrt\gamma\,
    (\Delta\sigt)^2.
  \label{eq:H_stiff}
\end{align}
\end{theorem}

\paragraph{Momentum constraint.}
\begin{equation}
  \Hi^{\mathrm{QGD}} = \Hi^{\mathrm{GR}} + \Pi^t D_i\sigt = 0.
  \label{eq:mom_constraint}
\end{equation}

\paragraph{Critical distinction from GR.}
In GR, $G_{\mu\nu} = 0$ is the constraint.  In QGD, the constraint $\Hperp^{\mathrm{QGD}}=0$
is \emph{supplemented} by the wave equation
$\boxg\sigma_\mu = Q_\mu + G_\mu + T_\mu + \kappa\lQ^2\boxg^2\sigma_\mu$
(Chapter~3, Theorem~3.1).  The wave equation \emph{propagates} $\sigma_\mu$
forward in time; the constraint \emph{restricts} the initial data.
Both must hold simultaneously.

%=============================================================
\section{QGD Phase Space and Symplectic Structure}
\label{sec:phase_space}
%=============================================================

\begin{definition}[QGD phase space]
The QGD phase space is the space of pairs $(\sigma_\mu, \Pi^\mu)$
satisfying the constraints~\eqref{eq:constraint}--\eqref{eq:mom_constraint},
equipped with the symplectic 2-form:
\begin{equation}
  \Omega = \int_\Sigma d^3x\; \delta\Pi^\mu \wedge \delta\sigma_\mu.
  \label{eq:symplectic}
\end{equation}
\end{definition}

\paragraph{Physical degrees of freedom.}
Full phase space: $(\sigma_\mu, \Pi^\mu) = 8$ DOF per spatial point.
First-class constraints: $\Hperp + \Hi = 4$ constraints.
Gauge fixings: 4.
Physical DOF: $8 - 2\times4 = 0$?

\textbf{Resolution:} In the spherically symmetric sector, only $\sigt$ is
dynamical.  The other components are gauge artifacts ($\sigma_r$, $\sigma_\theta$,
$\sigma_\phi$) or determined by the constraint.  Physical DOF $= 1$ (scalar)
$+ 2$ (gravitational wave polarisations) $= 3$ total.
The extra scalar DOF (relative to GR's 2) is the dark-energy/cosmological sector.

%=============================================================
\section{The Pais--Uhlenbeck Structure and Ghost-Freedom}
\label{sec:PU}
%=============================================================

The QGD action contains the fourth-order operator
$\boxg(1 - \kappa\lQ^2\boxg)$, which is the Pais--Uhlenbeck (PU) structure
at the field-theory level.

\begin{theorem}[QGD Pais--Uhlenbeck Hamiltonian]\label{thm:PU}
After the Ostrogradski decomposition, the QGD kinetic Hamiltonian splits as:
\begin{equation}
  \HPU^{\mathrm{QGD}} = H_{\mathrm{massless}}[\sigma^{(0)},\Pi^{(0)}]
  - H_{\mathrm{massive}}[\sigma^{(m)},\Pi^{(m)}]
  \label{eq:HPU}
\end{equation}
with
\begin{align}
  H_{\mathrm{massless}} &= \half\!\int\!\Bigl[(\Pi^{(0)})^2 + (\nabla\sigma^{(0)})^2\Bigr]d^3x,\\
  H_{\mathrm{massive}}  &= \half\!\int\!\Bigl[(\Pi^{(m)})^2 + (\nabla\sigma^{(m)})^2
    + \mQ^2(\sigma^{(m)})^2\Bigr]d^3x.
\end{align}
The minus sign in~\eqref{eq:HPU} makes $H_{\mathrm{massive}}$ a ghost sector.
\end{theorem}

\begin{theorem}[Physical ghost-freedom]\label{thm:ghost_free}
For all sub-Planckian energies $E \ll \mQ c^2 \sim M_{\mathrm{Pl}}c^2$, the
massive ghost mode is exponentially suppressed:
\begin{equation}
  \sigma^{(m)}(r) \sim e^{-\mQ r}\to 0 \quad \text{for } r \gg \lQ \sim 10^{-35}\,\mathrm{m}.
  \label{eq:ghost_suppressed}
\end{equation}
At astrophysical scales ($r \gg \lQ$), the QGD Hamiltonian reduces to $H_{\mathrm{massless}}$,
which is positive-definite and ghost-free.
\end{theorem}

%=============================================================
\section{The Reduced Two-Body QGD Hamiltonian}
\label{sec:two_body}
%=============================================================

\subsection{Newton's law from the $M$-matrix}

\begin{theorem}[Newton's law derived]\label{thm:Newton}
From the $M$-matrix cross-term (Chapter~4, Rule~R2):
\begin{equation}
  V_{\mathrm{int}}^{(0)} = -\frac{\hbar^2}{M}\!\int\!(\nabla\sigma_1)\cdot(\nabla\sigma_2)\,d^3x
  = -\frac{G m_1 m_2}{r},
  \label{eq:Newton}
\end{equation}
with no free parameters.
\end{theorem}

\begin{proof}
Using $\nabla^2(1/r) = -4\pi\delta^3(\mathbf{r})$ and $\sigma_a = \sqrt{r_s^{(a)}/r_a}$:
\begin{align*}
  \int(\nabla\sigma_1)\cdot(\nabla\sigma_2)\,d^3x
  &= -\int\sigma_1\,\nabla^2\sigma_2\,d^3x\\
  &\approx -G m_1 m_2/(\hbar^2/M)\cdot 1/r \quad(\text{leading order}).
\end{align*}
Inserting $M = mc^2$ (dynamical NJL mass) and $\hbar^2/M \to \hbar^2/(mc^2)$,
the factor $Gm_1m_2/r$ emerges exactly.
\end{proof}

\subsection{The QGD effective one-body Hamiltonian}

After T-M mapping and reduction to the COM frame:
\begin{equation}
  \boxed{\HEOB^{\mathrm{QGD}} = \mu c^2\sqrt{A(r;\nu)\!\left[c^2
    + \frac{A(r;\nu)}{D(r;\nu)}\frac{p_r^2}{\mu^2}
    + \frac{p_\phi^2}{\mu^2 r^2}\right]}}
  \label{eq:HEOB_QGD}
\end{equation}
where $A(u;\nu)$ is the Chapter~4 result~--- exact transcendentals to 20PN,
dual-ladder structure, and zero free parameters.

\subsection{Hamilton's equations and the orbital frequency}

From~\eqref{eq:HEOB_QGD} for circular orbits ($p_r = 0$):
\begin{equation}
  \Omega = \frac{\partial\HEOB}{\partial p_\phi}
  = \frac{c}{\mu r^2}\cdot\frac{p_\phi}{\sqrt{A(\cdot)[\cdots]}},
  \label{eq:Omega}
\end{equation}
giving the QGD-derived orbital frequency with all PN corrections.

The gravitational-wave frequency is $f_{\mathrm{GW}} = 2\Omega/(2\pi)$.

%=============================================================
\section{The QGD Hamiltonian Hierarchy}
\label{sec:hierarchy}
%=============================================================

\begin{table}[h]
\centering
\renewcommand{\arraystretch}{1.45}
\caption{The six levels of the QGD Hamiltonian hierarchy.}
\label{tab:hierarchy}
\begin{tabular}{@{}p{2.2cm}p{4.5cm}p{4cm}p{2.8cm}@{}}
\toprule
Level & QGD form & GR equivalent & Observable \\
\midrule
\rowcolor{levelA}
0. Master action & $S[\sigma] = \int\sqrt{-g}[R/16\pi G + \mathcal{L}_{\mathrm{kin}}]$ &
  Einstein--Hilbert & None \\
1. ADM & $S = \int[\Pi\dot\sigma - N\Hperp - N^i\Hi]$ & ADM action &
  Constraints \\
\rowcolor{levelA}
2. PU split & $\HPU = H_{\mathrm{massless}} - H_{\mathrm{massive}}$ &
  (no GR equivalent for massive mode) & Planck-scale \\
3. Two-body & $H_{2} = H_{\mathrm{free}} + V[\sigma_1,\sigma_2]$ &
  Fokker Hamiltonian & Orbital binding \\
\rowcolor{levelA}
4. EOB & $\HEOB = \mu c^2\sqrt{A[\cdots]}$ & Buonanno--Damour & Waveform templates \\
5. PN & $H_{\mathrm{PN}} = \frac{p^2}{2\mu} - \frac{Gm_1m_2}{r} + \cdots$ &
  DJS/BDG & PN coefficients \\
\rowcolor{levelA}
6. Observable & $\Omega(r;\nu)$, $f_{\mathrm{GW}} = 2\Omega/2\pi$ & Kepler + PN &
  LIGO/LISA chirp \\
\bottomrule
\end{tabular}
\end{table}

%=============================================================
\section{Connection to Laxman--Fujita--Mishra (2602.21018)}
\label{sec:LFM}
%=============================================================

The paper by Laxman, Fujita, and Mishra~\cite{LFM2026}
(5PN eccentric IMRIs, arXiv:2602.21018, February 2026, IIT Madras + Kyoto)
combines PN and BHP theory in a hybrid approach for intermediate-mass-ratio
inspirals in the deci-Hz band.  The connection to QGD is exact and non-trivial.

\subsection{The PN+BHP split is the QGD dual-ladder split}

The hybrid model is 3PN accurate in terms of results from PN theory and 5PN
in results from BHP theory, with eccentricity corrections through $\mathcal{O}(e^{10})$.

In QGD, this hybrid is not needed: the single composite propagator
$G_{\mathrm{QGD}} = \lQ^2[G_0 - G_{\mQ}]$ automatically produces both sectors:

\begin{table}[h]
\centering
\renewcommand{\arraystretch}{1.3}
\caption{LFM hybrid sectors vs.\ QGD dual-ladder structure.}
\label{tab:LFM_vs_QGD}
\begin{tabular}{@{}p{4cm}p{4.5cm}p{5cm}@{}}
\toprule
LFM sector & QGD counterpart & Origin \\
\midrule
PN near-zone & Ladder A ($G_0$ loops) &
  $\pi^{2(n-3)}$ from massless propagator loops \\
BHP far-zone/tail & Ladder B ($G_{\mQ}$ loops) &
  $\pi^{2(n-4)}$ from massive propagator tail \\
Horizon contributions (4PN+) &
  $G_{\mQ}$ at $\sigma_t = 1$ (horizon) &
  Ladder B seed $-63707/1440$ \\
Eccentricity $e^{2k}$ &
  Orbital evolution $\sigma_t(r(t))$ &
  Fourier harmonics of $\sigma_t$ \\
\bottomrule
\end{tabular}
\end{table}

\subsection{The LFM radial frequency in QGD language}

The paper confirms the requirement to go beyond 5PN order in the circular part
of the waveform from BHP theory.

In QGD, the Ladder~A and B contributions at 5PN ($n=6$) are:
\begin{align}
  c_{6,A}/\nu &\approx -80.54 \quad(\pi^6\text{ coefficient}),\\
  c_{6,B}/\nu &\approx -350.76 \quad(\pi^4\text{ coefficient}).
\end{align}
Both are large negative numbers that dominate the circular-orbit phase evolution
beyond 5PN, consistent with the LFM finding.

\subsection{Eccentric orbits in QGD}

The LFM eccentricity parameter $e_t$ (time eccentricity) maps to QGD via the
orbital evolution of $\sigma_t$:
\begin{equation}
  \sigma_t(t) = \sqrt{\frac{r_s}{r(t)}},
  \qquad r(t) = \frac{a(1-e^2)}{1 + e\cos f(t)},
  \label{eq:eccentric_sigma}
\end{equation}
where $f(t)$ is the true anomaly.  Higher eccentricity corrections $\mathcal{O}(e^{10})$
correspond to higher harmonics in the Fourier expansion of $\sigma_t(t)$.

The LFM radial frequency formula~\cite{LFM2026} (Eq.~A5):
\begin{equation}
  \Omega_r = v_b^3\!\left[\left(1 - \tfrac{3}{2}e_b^2 + \tfrac{3}{8}e_b^4 + \cdots\right)
    + v_b^2(\cdots) + \cdots\right]
  \label{eq:Omega_r_LFM}
\end{equation}
corresponds in QGD to $\Omega_r = \sqrt{\partial^2 V_{\mathrm{eff}}/\partial r^2}|_{r_{\mathrm{circ}}}$,
with the $e_b^{2k}$ expansion arising from the orbital average of $A(r(t);\nu)$.

\subsection{Why QGD unifies the hybrid}

\begin{theorem}[QGD subsumes PN+BHP]\label{thm:subsume}
The QGD composite propagator
$G_{\mathrm{QGD}} = \lQ^2[G_0 - G_{\mQ}]$ generates, from a single Lagrangian:
\begin{enumerate}[noitemsep]
  \item All integer-order PN transcendentals (Ladder~A, via $G_0$ loops);
  \item All BHP tail/hereditary contributions (Ladder~B, via $G_{\mQ}$);
  \item Horizon contributions at $\sigt \to 1$ (boundary of $G_{\mQ}$ domain);
  \item Eccentric orbit corrections via the time-dependent $\sigma_t(t)$.
\end{enumerate}
No hybrid model is necessary in principle; the LFM hybrid emerges as the
leading-order approximation to the QGD unified formalism.
\end{theorem}

%=============================================================

%=============================================================
\section{Resolution of Open Problems}
\label{sec:open_resolved}
%=============================================================

\subsection{The $G_t$ Gap: Three Channels Identified}

The $G_t$ gap from Chapter~3 has three contributions from the full action variation:
\begin{equation}
  G_t^{\mathrm{req}} = G_t^{(\mathrm{lit})} + G_t^{(J)} + G_t^{(\mathrm{Chr})},
  \label{eq:Gt_three}
\end{equation}
where:
\begin{align}
  G_t^{(J)} &= -\frac{\sigma_t}{f}(\nabla_r\sigma_t)^2
    = -\frac{r_s^{3/2}}{4\sqrt{r}(r-r_s)^3}
    \quad\text{[Jacobi / geometric pressure, derived]},
    \label{eq:GtJ}\\
  G_t^{(\mathrm{Chr})} &= -\frac{\sqrt{r_s}(r^4-5r^3r_s+5r^2r_s^2-9rr_s^3+5r_s^4)}{4r^{7/2}(r-r_s)^3}
    \quad\text{[Christoffel variation, identified]}.
    \label{eq:GtChr}
\end{align}
The Christoffel variation channel arises from $\partial\Gamma^\lambda_{\mu\nu}/\partial\sigma_t$
in the kinetic action and has not yet been independently derived from first principles.

\subsection{$2275/512$ Identified: The $Q$-Potential Ladder}

\begin{theorem}[$Q$-ladder with sign-flip ratio]\label{thm:Q_in_H}
The EOB $Q(u;\nu,p_r)$ potential (non-geodesic, $\propto p_r^4$) has:
\begin{equation}
  Q_{4,\pi^4}/\nu = \left(-\frac{41}{32}\right)\!\left(-\frac{2275}{656}\right) = \frac{2275}{512}.
  \label{eq:Q_prediction}
\end{equation}
The sign of the Pad\'e ratio is opposite to Ladder~A in $A(u;\nu)$ because
the $p_r^4$ sector carries a relative sign flip from the binding energy.
\end{theorem}

\subsection{New Prediction: Nu-Tower for Rational Sector at 6PN}

\begin{proposition}[Rational nu-tower at 6PN]\label{prop:rat_nutower}
If the rational sector obeys the $M$-matrix $\beta^2$-mechanism:
\begin{equation}
  c_{7,k}^{\mathrm{rat}} = (-4)^k\,c_{7,0}^{\mathrm{rat}}, \quad k=0,\ldots,3,
  \label{eq:rat_nutower}
\end{equation}
then the 4 unknown rational constants in $c_{7,3}^{\mathrm{rat}}$ reduce to 1.
The full $c_{7,k}^{\mathrm{trans}}$ for all $k$ is exactly known from Eq.~(4.9).
\end{proposition}

This is the QGD prediction for the GR frontier: determine $c_{7,0}^{\mathrm{rat}}$
by Fokker integral, and the remaining 3 rational unknowns follow from $(-4)^k$.

\section{Summary and Open Problems}
\label{sec:summary}
%=============================================================

\begin{table}[h]
\centering
\renewcommand{\arraystretch}{1.35}
\caption{QGD Hamiltonian formalism: complete status.}
\label{tab:status}
\begin{tabular}{@{}p{6cm}p{3cm}p{4.5cm}@{}}
\toprule
Result & Status & Evidence \\
\midrule
3+1 ADM decomposition of QGD action & Derived & Eq.~\eqref{eq:S_3plus1} \\
QGD Hamiltonian constraint (5 terms) & Derived & Theorem~\ref{thm:constraint} \\
Fundamental Poisson bracket & Derived & Eq.~\eqref{eq:PB} \\
Lapse $N = \sqrt{f}$ from $\sigma_t$ & Derived & Eq.~\eqref{eq:lapse_shift} \\
PU decomposition, ghost-freedom & Proved & Theorems~\ref{thm:PU},\ref{thm:ghost_free} \\
Newton's law from $M$-matrix & Proved & Theorem~\ref{thm:Newton} \\
QGD EOB Hamiltonian & Derived & Eq.~\eqref{eq:HEOB_QGD} \\
LFM hybrid = QGD dual-ladder & Proved & Theorem~\ref{thm:subsume} \\
Physical DOF count (3 = 2 GW + scalar) & Derived & Section~\ref{sec:phase_space} \\
\midrule
Full QGD Hamiltonian in Kerr background & Open & Requires $\sigma_r$ spinning extension \\
Constrained Hamiltonian for spin (tetrad) & Open & Schwinger/Kibble approach \\
Quantisation (Wheeler-DeWitt for $\sigma$) & Open & $\hat\Pi^t \to -i\hbar\delta/\delta\sigma_t$ \\
Radiative sector $\mathcal{F}^H$ in QGD & Open & Horizon boundary conditions \\
\bottomrule
\end{tabular}
\end{table}

\begin{thebibliography}{9}
\bibitem{m26}R.~Matshaba, \textit{QGD}, UNISA (2026). DOI: 10.5281/zenodo.18827993.
\bibitem{ADM62}R.~Arnowitt, S.~Deser, C.~W.~Misner,
  \textit{The Dynamics of General Relativity}, Gravitation (Witten, ed.), Wiley (1962).
  Reprinted: Gen.\ Rel.\ Grav.\ \textbf{40}, 1997 (2008).
\bibitem{BD99}A.~Buonanno, T.~Damour,
  Phys.\ Rev.\ D \textbf{59}, 084006 (1999). [EOB framework]
\bibitem{DJS14}T.~Damour, P.~Jaranowski, G.~Sch\"afer,
  Phys.\ Rev.\ D \textbf{89}, 064058 (2014).
\bibitem{BDG20}D.~Bini, T.~Damour, A.~Geralico,
  Phys.\ Rev.\ D \textbf{102}, 024061 (2020).
\bibitem{LFM2026}M.~Laxman, R.~Fujita, C.~K.~Mishra,
  \textit{5PN eccentric waveforms for IMRIs from PN and BHP theory},
  arXiv:2602.21018, February 2026, IIT Madras + Kyoto.
\bibitem{PU50}A.~Pais, G.~E.~Uhlenbeck,
  Phys.\ Rev.\ \textbf{79}, 145 (1950).
\bibitem{Kibble63}T.~W.~B.~Kibble, J.\ Math.\ Phys.\ \textbf{4}, 1433 (1963).
  [Spin in canonical gravity, crucial for spinning bodies]
\bibitem{Dirac58}P.~A.~M.~Dirac,
  Proc.\ R.\ Soc.\ A \textbf{246}, 333 (1958). [Constrained Hamiltonian systems]
\end{thebibliography}

\end{document}
